\begin{lemma}
\label{lemma-projective-modulo-two-ideals}
Let $R$ be a ring. Let $I, J \subset R$ be ideals such that $I \cap J = 0$.
Let $P$ be an $R$-module such that $P/IP$ is a projective $R/I$-module and
$P/JP$ is a projective $R/J$-module. Then $P$ is a projective $R$-module.
\end{lemma}

\begin{proof}
Choose a surjection $p : F \to P$ where $F$ is a free $R$-module.
Since $P/IP$ is a projective $R/I$-module, we can choose a map
$f : P/IP \to F/IF$ which is a right inverse to $p \bmod I$.
Consider the map $q : F/JF \to F/(I + J)F \times_{P/(I + J)P} P/JP$
induced by the quotient map $F/JF \to F/(I + J)F$ and $p$.
Note that $q$ is a surjective map of $R/J$-modules (small detail omitted).
Consider the map $f' : P/JP \to F/(I + J)F \times_{P/(I + J)P} P/JP$
induced by $f$ and the identity map.
Since $P/JP$ is a projective $R/J$-module and $q$ surjective,
we can choose a map $g : P/JP \to F/JF$ such that $q \circ g = f'$.
Then $f \bmod I + J = g \bmod I + J$. Since $F = F/JF \times_{F/(I + J)F} F/IF$
because $I \cap J = 0$, we conclude that we obtain a well defined homomorphism
$h = (f, g) : P \to F$ of $R$-modules. The map $a = p \circ h : P \to P$
reduces to the identity modulo $I$ and modulo $J$. Then $a$ also induces
the identity map on the submodule $IP$: if $x = \sum i_\alpha x_\alpha$
with $i_\alpha \in I$ and $x_\alpha \in P$, then
$a(x) = \sum i_\alpha a(x_\alpha) = \sum i_\alpha x_\alpha = x$
because $a(x_\alpha) = x_\alpha + j_\alpha$ with $j_\alpha \in J$
and $i_\alpha j_\alpha = 0$. Thus $a : P \to P$ acts as the identity
on $IP$ and on $P/IP$ and we conclude that $a$ is an automorphism of $P$
(small detail omitted). Thus $P$ is a summand of $F$, whence projective.
\end{proof}
